\setcounter{demsode}{2} %%% THỨ TỰ ĐỀ THẦY - CÔ ĐÃ NHẬN
\begin{name}
	{Biên soạn: Phạm Thế Sinh\\
	Phản biện: Tư Đô Nguyên}
	{ÔN TẬP HỌC KỲ II KHỐI 10 NĂM 2024}
\end{name}

\cautn
\Opensolutionfile{ans}[ans/0-CTST-TN-SO-2]
\setcounter{ex}{0}

%%% CÂU 1
\begin{ex} %[DA-HK2-K10,11-2024 - PT Sinh] %[0D8N2-2]
	Có bao nhiêu cách sắp xếp chỗ ngồi cho $5$ học sinh vào $5$ ghế xếp thành một dãy?
	\choice
	{$90$}
	{$60$}
	{\True $120$}
	{$240$}
	\loigiai{
	}
\end{ex}

%%% CÂU 2
\begin{ex}%[DA-HK2-K10,11-2024 - Phạm Thế Sinh] %[0D0N1-2]
	Cho phép thử có không gian mẫu $\Omega=\{1;2;3;4;5\}$. Tìm cặp biến cố không đối nhau trong các cặp biến cố sau?
	\choice
	{\True $C=\{1;4\}$ và $D=\{2;3\}$}
	{$E=\{1;3;5\}$ và $F=\{2;4\}$}
	{$A=\{1\}$ và $B=\{2;3;4;5\}$}
	{$\Omega$ và $\varnothing$}
	\loigiai{
	}
\end{ex}

%%==========Câu 3
\begin{ex}%[DA-HK2-K10,11-2024 - Phạm Thế Sinh] %[0H9H3-2]
Đường thẳng $\Delta$ đi qua điểm $M(-1;4)$ và có vectơ pháp tuyến $\vec{n}=(2;3)$ có phương trình tổng quát là
\choice
{$-x+4y+10=0$}
{\True $2x+3y-10=0$}
{$2x+3y+10=0$}
{$-x+4y-10=0$}
\loigiai{

}
\end{ex}

%%==========Câu 4
\begin{ex}%[DA-HK2-K10,11-2024 - Phạm Thế Sinh] %[0D8N2-5]
Từ một bó hoa hồng gồm $3$ bông hồng trắng, $5$ bông hồng đỏ và $6$ bông hồng vàng, có bao nhiêu cách chọn ra một bông hồng?
\choice
{$8$}
{$11$}
{\True $14$}
{$90$}
\loigiai{

}
\end{ex}

%%==========Câu 5
\begin{ex}%[DA-HK2-K10,11-2024 - Phạm Thế Sinh] %[0D0N2-4]
Một nhóm học sinh gồm $10$ học sinh nam và $5$ học sinh nữ. Giáo viên chọn ngẫu nhiên một học sinh đi lên bảng làm bài tập. Tính xác suất chọn được một học sinh nữ?
\choice
{$\dfrac{1}{2}$}
{$\dfrac{1}{10}$}
{$\dfrac{1}{5}$}
{\True $\dfrac{1}{3}$}
\loigiai{

}
\end{ex}

%%==========Câu 6
\begin{ex}%[DA-HK2-K10,11-2024 - Phạm Thế Sinh] %[0H9N3-1]
Cho đường thẳng $\Delta\colon 2x-y+1=0$. Điểm nào sau đây nằm trên đường thẳng $\Delta$?
\choice
{$C\left(\dfrac{1}{2};-2\right)$}
{$D(0;-1)$}
{$A(1;1)$}
{\True $B\left(\dfrac{1}{2};2\right)$}
\loigiai{

}
\end{ex}

%%==========Câu 7
\begin{ex}%[DA-HK2-K10,11-2024 - Phạm Thế Sinh] %[0H9N4-1]
Xác định tâm của đường tròn có phương trình $2x^2+2y^2-8x+4y-1=0$
\choice
{$(2;-1)$}
{$(8;-4)$}
{$(-8;4)$}
{\True $(4;-2)$}
\loigiai{

}
\end{ex}

%%==========Câu 8
\begin{ex}%[DA-HK2-K10,11-2024 - Phạm Thế Sinh] %[0D7H2-1]
Tập nghiệm của bất phương trình $-x^2+2x+7<0$
\choice
{$S=(-\infty;1-2 \sqrt{2}) \cup(1+2\sqrt{2};+\infty)$}
{\True $S=(1-2\sqrt{2};1+2 \sqrt{2})$}
{$S=\varnothing$}
{$S=\mathbb{R}$}
\loigiai{

}
\end{ex}

%%==========Câu 9
\begin{ex}%[DA-HK2-K10,11-2024 - Phạm Thế Sinh] %[0H9N5-2]
Phương trình nào dưới đây có đồ thị là một hình elip?
\choice
{$\dfrac{x^2}{25}+\dfrac{y^3}{9}=1$}
{\True $\dfrac{x^2}{9}+\dfrac{y^2}{4}=2$}
{$\dfrac{x}{18}+\dfrac{y^2}{9}=2$}
{$\dfrac{x}{16}+\dfrac{y}{9}=1$}
\loigiai{

}
\end{ex}

%%==========Câu 10
\begin{ex}%[DA-HK2-K10,11-2024 - Phạm Thế Sinh] %[0D7H3-4]
Tập nghiệm của phương trình $(x^2-x-2) \sqrt{x-1}=0$ là
\choice
{$\{1;2\}$}
{\True $\{-1;1;2\}$}
{$[1;2]$}
{$\{-1;2\}$}
\loigiai{

}
\end{ex}

%%==========Câu 11
\begin{ex}%[DA-HK2-K10,11-2024 - Phạm Thế Sinh] %[0D8N1-2]
Từ nhà bạn An đến nhà bạn Bình có $3$ con đường đi, từ nhà bạn Bình đến nhà bạn Cường có $2$ con đường đi. Hỏi có bao nhiêu cách chọn đường đi từ nhà bạn An đến nhà bạn Cường và phải đi qua nhà bạn Bình?
\choice
{$2$}
{\True $6$}
{$5$}
{$3$}
\loigiai{

}
\end{ex}

%%==========Câu 12
\begin{ex}%[DA-HK2-K10,11-2024 - Phạm Thế Sinh] %[0D7N1-1]
Tìm khẳng định đúng trong các khẳng định sau?
\choice
{$f(x)=3 x^3+2x-1$ là tam thức bậc hai}
{$f(x)=x^4-x^2+1$ là tam thức bậc hai}
{\True $f(x)=3x^2+2x-5$ là tam thức bậc hai}
{$f(x)=2x-4$ là tam thức bậc hai}
\loigiai{

}
\end{ex}


\Closesolutionfile{ans}
\indapan{6}{ans/0-CTST-TN-SO-2}

\cauds
\Opensolutionfile{ans}[ans/0-CTST-DS-SO-2]
\setcounter{ex}{0}
Học sinh trả lời từ câu $1$ đến câu $4$. Trong mỗi ý a), b), c), d) ở mỗi câu, học sinh chọn đúng hoặc sai.
%%% CÂU 1
\begin{ex} %[DA-HK2-K10,11-2024 - Phạm Thế Sinh] %[0D0V2-5]
Một hộp đựng $9$ thẻ được đánh số từ $1$ đến $9$. Rút ngẫu nhiên $5$ thẻ. Khi đó:
	\choiceTF
	{\True Xác suất ``Các thẻ ghi số $1$, $2$, $3$ được rút'' bằng $\dfrac{5}{42}$}
	{Xác suất ``Có đúng $1$ trong $3$ thẻ ghi sỗ $1$, $2$, $3$ được rút'' bằng $\dfrac{6}{11}$} %5/14
	{\True Xác suất ``Không thẻ nào trong $3$ thẻ ghi số $1$, $2$, $3$ được rút'' bằng $\dfrac{1}{21}$}
	{\True Xác suất ``Có ít nhất một trong $3$ thẻ ghi số $1$, $2$, $3$ được rút'' bằng $\dfrac{20}{21}$}
	\loigiai{
		
	}
\end{ex}

%%% CÂU 2
\begin{ex}%[DA-HK2-K10,11-2024 - Phạm Thế Sinh] %[0D0V2-4]
Một nhóm gồm $8$ nam và $7$ nữ. Chọn ngẫu nhiên $5$ bạn. Khi đó:
	\choiceTF
	{\True Số phần tử của không gian mẫu: $3003$}
	{\True Xác suất để có đúng $1$ bạn nữ bằng $\dfrac{70}{429}$}
	{\True Xác suất để có $3$ nam và $2$ nữ bằng $\dfrac{56}{143}$}
	{Xác suất để có cả nam lẫn nữ mà nam nhiều hơn nữ bằng $\dfrac{23}{429}$} %60/429
	\loigiai{
	
	}
\end{ex}

%%% CÂU 3
\begin{ex}%[DA-HK2-K10,11-2024 - Phạm Thế Sinh] %[0D7V3-1]
Cho phương trình $(x+1)\left(\sqrt{x+4}-\sqrt{-x^2+4 x+14}\right)=0$ (*). Khi đó:
	\choiceTF
	{Điều kiện: $x \geq 4$}
	{\True Phương trình (*) có $3$ nghiệm phân biệt}
	{Các nghiệm của phương trình (*) nhỏ hơn $5$}
	{\True Tổng các nghiệm của phương trình (*) bằng $2$}
	\loigiai{
		
	}
\end{ex}

%%% CÂU 4
\begin{ex}%[DA-HK2-K10,11-2024 - Phạm Thế Sinh] %[0H9C3-7]
Cho tam giác $ABC$ có phương trình của đường thẳng $BC$ là $7x+5y-8=0$, phương trình các đường cao kẻ từ $B$, $C$ lần lượt là $9x-3y-4=0$, $x+y-2=0$. Khi đó:
	\choiceTF
	{\True Điểm $B$ có tọa độ là $\left(\dfrac{2}{3};\dfrac{2}{3}\right)$}
	{\True Điểm $C$ có tọa độ là $(-1;3)$}
	{Phương trình đường cao kẻ từ $A$ là $5x-7y-6=0$} %5x-7y+4=0
	{Phương trình đường trung tuyến kẻ từ $A$ là $x-13y+4=0$}
	\loigiai{
	
	}
\end{ex}
\Closesolutionfile{ans}
\indapan{4}{ans/0-CTST-DS-SO-2}

\caukq
\Opensolutionfile{ans}[ans/0-CTST-KQ-SO-2]
\setcounter{ex}{0}
Học sinh trả lời từ câu $1$ đến câu $6$.
%%% CÂU 1
\begin{ex}%[DA-HK2-K10,11-2024 - Phạm Thế Sinh] %[0D0H2-5]
Chọn ngẫu nhiên $2$ số trong tập hợp $X=\{1;2;3;\cdots ; 50\}$. Tính xác suất của biến cố $A$: ``Hai số được chọn là số chẵn'' (làm tròn đến chữ số phần trăm).
	\shortans[]{$0{,}24$}
	\loigiai{
	}
\end{ex}

%%% CÂU 2
\begin{ex}%[DA-HK2-K10,11-2024 - Phạm Thế Sinh] %[0D0H2-4]
Một lớp học có $26$ bạn nam và $20$ bạn nữ. Chọn ngẫu nhiên một bạn trong lớp. Tính xác suất để bạn được chọn là nam (làm tròn đến chữ số phần trăm).
	\shortans[]{$0{,}57$}
	\loigiai{
	}
\end{ex}

%%% CÂU 3
\begin{ex}%[DA-HK2-K10,11-2024 - Phạm Thế Sinh] %[0H9H4-2]
Trong mặt phẳng tọa độ $Oxy$, tính bán kính đường tròn qua $A(2;-1)$ và tiếp xúc với $Ox$.

	\shortans[]{$2$}
	\loigiai{
	}
\end{ex}

%%% CÂU 4
\begin{ex}%[DA-HK2-K10,11-2024 - Phạm Thế Sinh] %[0H9V4-2]
Trong mặt phẳng tọa độ $Oxy$, tính bán kính đường tròn đi qua hai điểm $A$, $B$ và có tâm $I$ nằm trên đường thẳng $\Delta$, với $A(0;4)$, $B(2;6)$, $\Delta\colon x-2y+5=0$ (làm tròn đến chữ số phần trăm).
	\shortans[]{$2{,}36$}
	
	\loigiai{
	}
\end{ex}

%%% CÂU 5
\begin{ex}%[DA-HK2-K10,11-2024 - Phạm Thế Sinh] %[0H9V4-7]
Một cái cổng hình bán nguyệt rộng $8{,}4$ m, cao $4{,}2$ m như hình vẽ. Một đường dưới cổng được chia làm hai làn cho xe ra vào. Một chiếc xe tải rộng $2{,}2$ m, cao $2{,}6$ m đi đúng làn đường quy định, hỏi mép thùng xe phía trên còn cách trần cổng bao nhiêu mét? (tính theo phương thẳng đứng và làm tròn đến hàng phần mười)
\begin{center}
\begin{tikzpicture}[>=stealth,font=\footnotesize,scale=1,declare function={R=2;r=1.6;l=1.4;a=sqrt(r^2-(l-R+r)^2);}]
	\path (0,0) coordinate (O) (0,R-r) coordinate (O') (-R,0) coordinate (A) 
	(R,0) coordinate (B)(-R,R) coordinate (C) (R,R) coordinate (D);
	\draw[<->,dashed] (O)--node[right,pos=.4]{$4{,}2$ m}(90:R);
	\draw[<->,dashed] (A)--node[below]{$8{,}4$ m}(B);
	\draw ($(O')+(0:r)$) arc (0:180:r) (B) arc (0:180:R) (-R+l,l)--++(-135:3.5) (R-l,l)--++(-45:3.5);
	\draw[double,double distance = 5pt,shift={(-1pt,2pt)}] (-R,R) arc (136:44:{R*sqrt(2)});
	\draw[dashed] (-a,l)--(a,l);
	\fill[pattern=bricks] (-2.5*R,0)--(A)--(C)--+(180:1.5*R)--cycle (2.5*R,0)--(B)--(D)--+(0:1.5*R)--cycle;
	\draw (-2.5*R,0)--(A)--(C)--+(180:1.5*R) (2.5*R,0)--(B)--(D)--+(0:1.5*R);
\end{tikzpicture}
\end{center}
	\shortans[]{$1{,}0$}
	\loigiai{
	}
\end{ex}

%%% CÂU 6
\begin{ex}%[DA-HK2-K10,11-2024 - Phạm Thế Sinh] %[0H9V4-7]
Hình vẽ bên dưới mô phỏng một trạm thu phát sóng điện thoại di động đặt ở vị trí $I$ có tọa độ $(-2;1)$ trong mặt phẳng tọa độ (đơn vị trên hai trục là ki-lô-mét). Tính theo đường chim bay, xác định khoảng cách ngắn nhất để một người ở vị trí có tọa độ $(-3;4)$ di chuyển được tới vùng phủ sóng theo đơn vị ki-lô-mét (làm tròn kết quả đến hàng phần trăm). Biết rằng trạm thu phát sóng đó được thiết kế với bán kính phủ sóng $3\mathrm{~km}$.
\begin{center}
\begin{tikzpicture}[line join=round,font=\footnotesize,scale=.7]
	\path (0,0) coordinate (O) (-3,4) coordinate (M) (-2,1) coordinate (I) 
	node[shift={(120:1.2)},text width=2cm,align=center]{Trạm phát sóng};
	\draw[-stealth] (-5.5,0)--(2.5,0) node[below] {$x$};
	\draw[-stealth] (0,-2)--(0,4.5) node[left] {$y$};
	\draw (I) circle (3);
	\draw[dashed] (-2,0)node[below]{$-2$}|-(0,1) node[right]{$1$};
	\foreach \i/\j in {I/135,O/-135,M/90}
	\fill (\i) circle (1.3pt) node[shift={(\j:3mm)}]{$\i$};
\end{tikzpicture}
\end{center}
	\shortans[]{$0{,}16$}
	\loigiai{
	}
\end{ex}
\Closesolutionfile{ans}
\indapan{6}{ans/0-CTST-KQ-SO-2}